内部有效性最强的是同进程配对干预；版本、训练状态或接口同时变化时最弱。
\Cref{tab:8-threats}列出可反转各解释的因素。第3章的共享输入、明确分数契约、
计时数组和版本清单构成主要控制。重复调用估计条件计时波动，不估计独立重训
索引的分布；五次补证构建量化了固定工作点的recall变化，而非整条重训前沿分布；精度敏感性在
同一缓存模型下改变add／查询算术。

\par\medskip\noindent\begin{minipage}{\linewidth}
\centering\captionsetup{hypcap=false}
\captionof{table}[有效性威胁]{可能反转或缩小结论的威胁。}
\label{tab:8-threats}
\footnotesize\setlength{\tabcolsep}{4pt}
\begin{tabularx}{\linewidth}{@{}*{3}{>{\raggedright\arraybackslash}X}@{}}
\toprule
受影响主张 & 替代解释或反转 & 已有控制与剩余验证 \\
\midrule
执行收益 & 分数、并列或索引改变产生表面提速 & 配对状态与输出检查；
T4 严格 CPU 排序及表形式数值诊断；覆盖有限，候选 dump 已不保留 \\
预算质量 & 有限池抽样或漂移掩盖有害查询 & 留出GT、完整轨迹、子集和核对及FULL枚举；
未来查询及任意子集不变性未测 \\
竞争差距 & 网格密度、慢调用或API开销改变排名 & 两种recall估计、
保留慢调用、int8及TF32/FP32敏感性；T5 已显示小batch反转，跨设备比较仍未测 \\
容量／构建 & 瞬时分配超过常驻样本 & 阶段字节、重复空缓存构建及20\,ms设备峰值；
更短暂分配与完整重训前沿未测 \\
\bottomrule
\end{tabularx}
\end{minipage}\par\medskip

外部有效性限于一张 GPU、平方 L2 top-10、六份准备后的嵌入数据、E1 的 batch 1,000
及 T5 的 8/32。T5 高 recall 排名反转说明 E1 排名不能跨 batch 外推。
低目标比较必须保留已达标前沿端点；实测阈值结果揭示受网格影响的排名。
按位置移除查询不认证没有近重复嵌入，这些导出也不是带相关性标签的检索基准。
数据库更新、并发流量、部署延迟目标和变化的查询总体未测试。
压缩索引常驻结果不构成十亿规模或多设备容量证据。
这些局限给出把所观测工作区间迁移到部署前需要补的验证。
